% 第2章 完备市场理论

% 导言

\chapter{完备市场理论与宏观经济学中的均衡}

目前为止，我们都是从\textbf{计划者}（Social Planner）角度和\textbf{代表性家庭}（同质家庭）的假设下处理最优增长模型。
那么，如果市场结构是分散的（Decentralized），如何定义经济中的均衡？这样的市场结构下产生的均衡配置（allocation）有什么特征，和
计划者配置下的均衡有什么关系？代表性家庭的设定有什么含义？

理解上述问题离不开宏观经济学中的一个重要概念：完备市场（Complete Market）。
本章将先对完备市场理论进行介绍，引出两种市场交易结构和帕累托最优的概念（Jla、Edmond）；
接着，借助相关概念定义竞争均衡，并将其运用在最优增长模型上，研究竞争均衡下的福利问题；
最后，本章将介绍随机最优增长模型，在其上定义递归竞争均衡，引出后续对数值方法的进一步学习。



% 2.1 ==== Kruger Decentralized Equilibrium Lafo
\section{最优增长模型中的竞争均衡}

目前为止，我们都是从\textbf{计划者}（Social Planner）角度和\textbf{代表性家庭}（同质家庭）的假设下处理最优增长模型。
在后面我们将会证明，由此得到的配置（allocation）是帕累托有效（最优）的（Pareto Efficiency / Optimal）。从直觉上讲，
如果这一问题下的最优配置不是帕累托有效的，那么计划者的目标效用函数就存在提高空间。

为了更清晰的阐明宏观经济学中相关的均衡概念，现在先考虑如下问题：
当模型中包含企业（firm）和消费者（consumer）在市场上的互动时，经济中的竞争均衡（Competitive Equilibrium）是怎样的？
假设消费者拥有所有的生产要素（资本）并将其租赁给企业；假设代表性家庭拥有这些企业，索取其利润；假设产品市场和要素市场（劳动与资本）
是完全竞争的（perfectly competitive），这意味着家庭和企业都将价格视为给定（take as given）。


% 2.2 ==== 随机最优增长：引出不确定性的概念 Edmond
\section{随机最优增长模型}


% Edmond
\subsection{sequential approach using histories and contingent plans}


\subsection{recursive approach}


% 2.3 ==== 完备市场理论：抽象
\section{完备市场理论}


%
\subsection{计划者问题}


% AD介绍、与Pareto均衡的关系 （equilibrium allocation vs. planning allocation） Kruger、Edmond
\subsection{Arrow-Debreu Equilibrium}


%
\subsection{Sequential Equilibrium}


%
\subsection{两种均衡的等价性}


% Kruger 71
\subsection{Recursive Comptetitive Equilibrium}

在一些宏观模型中，均衡可能并非帕累托最优，也可能难以通过序列形式或Arrow-Debreu形式求解，这时我们通常会定义递归竞争均衡（Recursive Comptetitive Equilibrium）。



% 完备市场的意义：perfect risk sharing % consumption smoothing
\subsection{完备市场的意义}


% 
\subsection{RA设定原理}


% 2.4 ====
\section{完备市场理论的应用：资产定价}

